% Hexagon note, 23 Nov 2026. CC0 1.0.
% Mathematics: GPT-6 Astra (ChatGPT Pro mode). Text, corollary for analytically irreducible domains and the
% worked examples: Claude Opus 5.5 at effort high (editor).
\documentclass[11pt,a4paper]{article}
\usepackage[T1]{fontenc}
\usepackage{lmodern}
\usepackage{microtype}
\usepackage[margin=2.7cm]{geometry}
\usepackage{amsmath,amssymb,amsthm}
\PassOptionsToPackage{hyphens}{url}
\usepackage[hidelinks]{hyperref}
\emergencystretch=2em

\theoremstyle{plain}
\newtheorem{theorem}{Theorem}
\newtheorem{corollary}[theorem]{Corollary}
\newtheorem{lemma}[theorem]{Lemma}
\newtheorem{proposition}[theorem]{Proposition}
\theoremstyle{remark}
\newtheorem{example}[theorem]{Example}
\newtheorem{remark}[theorem]{Remark}

\newcommand{\F}{\mathbf F}
\newcommand{\Int}{\operatorname{Int}}
\newcommand{\Fn}{\mathcal F}
\newcommand{\workbenchurl}{https://github.com/KokunoYumeto/mathematics-workbenches/blob/main/works/WB-0004.md}

\title{Integer-valued polynomials over a one-dimensional analytically\\ irreducible local domain form a free module:\\ Elliott's non-flatness conjectures fail}
\author{}
\date{23 November 2026 \quad CC0 1.0}

\begin{document}
\maketitle

\begin{abstract}
\noindent
In 2011 J.~Elliott conjectured that the ring $\Int(D)$ of integer-valued polynomials is not flat over
$D=\F_2[[T^2,T^3]]$, nor over $D=\F_2+T\F_4[[T]]$, and more generally that some local, Noetherian,
one-dimensional, analytically irreducible domain has a non-flat $\Int(D)$. This spring H.~Ma, with the reasoning
system Rethlas, answered Problem~19 of Cahen, Fontana, Frisch and Glaz negatively, with a local domain whose
normalization has two maximal ideals. We prove that Elliott's conjectures fail. If $D$ is a subring of a
discrete valuation ring $V$ with finite residue field, and $D$ contains a nonzero ideal of $V$, then $\Int(D)$ is a
free $D$-module. Consequently $\Int(D)$ is free for every one-dimensional analytically irreducible Noetherian local
domain $D$. By Elliott's own theorems, the canonical maps $\Int(D)^{\otimes n}\to\Int(D^n)$ are then isomorphisms and
$\Int(D)$ carries a unique compatible biring structure. For Elliott's two rings we write down the polynomials
explicitly. The proof patches two free modules along the conductor square and uses a matrix identity of
Whitehead and a theorem of Kaplansky. The proof was produced by GPT-6 Astra (ChatGPT Pro mode). The corollary
for analytically irreducible domains and the explicit examples were written by Claude Opus 5.5 at effort high,
which also checked the proof; Claude Opus 5.5 at effort max checked the whole note again. No human has checked
this note.
\end{abstract}

\section{The question and the theorem}

Let $D$ be an integral domain with fraction field $K$. The ring of integer-valued polynomials is
$\Int(D)=\{f\in K[X]: f(D)\subseteq D\}$. It contains $D[X]$ and is in general larger, as $\binom{X}{2}\in\Int(\mathbf Z)$ shows. As a $D$-module, $\Int(D)$ is free when $D$ is
a Dedekind domain. It is locally free when $D$ is a Krull domain, or more generally a TV PvMD
\cite[Theorem~1.2]{E11a}. In the list of Cahen, Fontana, Frisch and Glaz, Elliott asks whether this persists
\cite[Problem~19, p.~9]{CFFG}: is $\Int(D)$ flat over every domain $D$ (19a)? Is it free (19b)?

Elliott expected the answer to be no, and named the rings where he expected it to fail. For a local domain $D$ whose
maximal ideal $M$ is finitely generated and not principal, $M^{-1}$ is an overring. If $\Int(D)$ is flat, then
$\Int(D,M^{-1})=M^{-1}\Int(D)$ \cite[Proposition~2.13, p.~12]{E11a}. For $D=\F_2+T\F_4[[T]]$, with $M^{-1}=\F_4[[T]]$,
and for $D=\F_2[[T^2,T^3]]$, with $M^{-1}=\F_2[[T]]$, he reduced non-flatness to the statement
\begin{equation}\label{eq:elliott}
 X^2+X\notin M\,\Int(D)
\end{equation}
\cite[Propositions~2.14 and~2.16, p.~12]{E11a}. He conjectured that~\eqref{eq:elliott} holds for both rings
\cite[Conjectures~2.15 and~2.17]{E11a}. Since the two rings are local, Noetherian, one-dimensional and
analytically irreducible, either conjecture would give his Conjecture~2.10 \cite[p.~11]{E11a}: \emph{some such domain
has a non-flat $\Int(D)$}. He recalled these conjectures as open in 2011 \cite[p.~2]{E11b}. In 2016 he still
conjectured that domains with non-free $\Int(D)$ exist \cite[p.~2]{E16b}.

\begin{theorem}\label{thm:main}
Let $V$ be a discrete valuation ring with finite residue field, uniformizer $\pi$ and fraction field $K$. Let
$R\subseteq V$ be a subring with the same unit, and suppose $\pi^mV\subseteq R$ for some $m\geq1$. Then $R$ is a
one-dimensional Noetherian local domain with fraction field $K$, and $\Int(R)$ is a free $R$-module of countably
infinite rank.
\end{theorem}

\begin{corollary}\label{cor:ai}
Let $D$ be a one-dimensional Noetherian local domain whose completion is a domain. Then $\Int(D)$ is a free
$D$-module. In particular Elliott's Conjectures~2.10, 2.15 and~2.17 of \cite{E11a} are false, and
$X^2+X\in M\,\Int(D)$ for both of his rings.
\end{corollary}

\begin{corollary}\label{cor:biring}
For every $D$ as in Corollary~\ref{cor:ai}, the canonical maps $\Int(D)^{\otimes_D n}\to\Int(D^n)$ are isomorphisms
for all $n$. There is a unique $D$-$D$-biring structure on $\Int(D)$ for which $D[X]\to\Int(D)$ is a biring
homomorphism.
\end{corollary}

In words: over a one-branch local domain of dimension one with finite normalization, integer-valued polynomials
behave exactly as over a discrete valuation ring, as far as module structure and the operations of composition and
tensor products are concerned. Corollary~\ref{cor:biring} follows from Theorem~\ref{thm:main} by theorems of Elliott:
if $\Int(D)$ is free, the canonical maps are isomorphisms \cite[Propositions~6.8(a), 6.10(d)]{E07}, as recalled in
\cite[p.~9]{CFFG} and \cite[p.~16]{E11b}, and then the biring exists and is unique \cite[p.~10]{CFFG}, by
\cite{E10}. We have not re-read the proofs in \cite{E07,E10}. Theorem~\ref{thm:main} gives Problems~19, 20 and~21
of \cite{CFFG} a positive answer on this class of domains.

The conjectured witness becomes explicit. For $D=\F_2[[t^2,t^3]]$ put
\[
 q=\frac{X^2+X}{t^2},\qquad h_2=\Bigl(\frac{q+q^2}{t}\Bigr)^4,\qquad h_1=q+t\,h_2 .
\]
Then $h_1,h_2\in\Int(D)$ and $X^2+X=t^2h_1+t^3h_2$, with $t^2,t^3\in M$ (Example~\ref{ex:cusp}). For
$D=\F_2+t\F_4[[t]]$, with $\F_4=\F_2(\omega)$, put $q=(X^2+X)/t$, $h_2=q+q^2$ and $h_1=q+\omega h_2$. Then
$X^2+X=t\,h_1+\omega t\,h_2$ (Example~\ref{ex:f4}).

\paragraph{Why this matters.}
Whether $\Int(D)$ is always flat had been open since Elliott's papers, and its negative answer this year
\cite{Ma} made it natural to ask where flatness fails. Ma's domain is $\F_2+t(t+1)T$ with $T=\F_2[t]$ localized at
the two primes $(t)$ and $(t+1)$ \cite[\S3]{Ma}. It is local and one-dimensional, but its normalization $T$ has two
maximal ideals: it has two branches. The same domain answers Problem~21 negatively \cite[Theorem~3.4]{JLSWXY}.
Theorem~\ref{thm:main} shows that one branch with finite normalization never fails. So the domains Elliott proposed
as the first examples of non-flatness are free, and the number of branches is the feature that separates the two
behaviours in dimension one.

\paragraph{The credit.}
The questions, the flatness criterion and the two candidate rings are Elliott's \cite{E11a,E11b,CFFG}. The first
non-flat example is Ma's, found with Rethlas \cite{Ma,Rethlas}. The proof of Theorem~\ref{thm:main} uses known tools,
cited where they are used:
\begin{itemize}
\item the Newton basis built from a sequence that minimizes products of differences, which is Bhargava's
$\pi$-ordering \cite{Bha};
\item the polynomial $(X^Q-X)/\pi$, which Elliott also uses \cite[Proposition~5.4, p.~16]{E11b};
\item Milnor's patching of modules over a conductor square \cite[\S2]{Mil};
\item Whitehead's lemma that $\operatorname{diag}(\alpha,\alpha^{-1})$ is a product of elementary matrices;
\item Kaplansky's theorem that projective modules over a local ring are free \cite{Kap}.
\end{itemize}
New here are Theorem~\ref{thm:main}, Corollaries~\ref{cor:ai} and~\ref{cor:biring} as consequences, and the explicit
identities above. We know of no earlier proof that $\Int(D)$ is flat for either of Elliott's two rings. The
standard monograph \cite{CC} was not available to us; Elliott, who cites it, still describes the question as open in
2011 and 2016 \cite{E11b,E16b}.

\section{Context}

For a Dedekind domain $\Int(D)$ is free \cite[Remark~II.3.7(iii)]{CC}, as cited in \cite[p.~2]{E11b}.
Elliott proved local freeness for Krull domains and TV PvMDs \cite[Theorem~1.2]{E11a}. Our domains are neither:
$\F_2[[t^2,t^3]]$ is not integrally closed, and its maximal ideal is $t$-maximal but not principal, which Elliott
notes places it outside his theorem \cite[p.~4]{E11a}.

Ma's theorem \cite[Theorem~1.1]{Ma} answers Problems~19a and~19b negatively. Its proof uses Elliott's criterion: the
polynomial $(X^2+X)/(t(t+1))$ lies in $\Int(D,T)$ but not in $T\,\Int(D)$ \cite[Proposition~4.1, pp.~4--5]{Ma}. The
argument was produced by Rethlas and checked by Ma \cite[\S6]{Ma}. Jiang, Li, Sun, Wang, Xiao and Yu use the same
domain to show that no compatible biring structure need exist, answering Problem~21 \cite[Theorem~3.4]{JLSWXY}. Both
results bear on this note in the same way. Their domain has two branches, and Theorem~\ref{thm:main} cannot apply,
because its normalization is not local. Our theorem covers every one-branch domain with finite normalization, and
none of theirs. Neither result implies the other. Together they show that, among one-dimensional local domains with
finite normalization, the branch count decides flatness and the existence of the biring.

\section{Proof overview}

The difficulty is that $\Int(R)$ is not finitely generated, so the usual Noetherian arguments do not apply: they
identify flat finitely generated modules with projective ones and check them locally. Elliott's criterion gives a
necessary condition for flatness, the identity $V\Int(R)=\Int(R,V)$. Even proving that identity directly means
writing every $V$-valued polynomial as a combination of $R$-valued ones. The examples above show that this is possible but
intricate.

The idea is to use the conductor square
\[
 R=V\times_{V/J}R/J,\qquad J=\pi^mV .
\]
It cuts $\Int(R)$ into two pieces that are easy to understand.
\begin{itemize}
\item Over the discrete valuation ring $V$, the module $B=\Int(R,V)$ is free, with an explicit Newton basis.
\item Modulo $J$, both $B$ and $\Int(R)$ become rings of locally constant functions on $R$, with values in the finite
rings $V/J$ and $R/J$. Such function modules are free, with bases obtained by refining congruence classes. The
digit polynomial $(X^Q-X)/\pi$ produces every congruence indicator, which shows that all such functions are values of
polynomials.
\end{itemize}
So $\Int(R)$ is patched from a free $V$-module and a free $R/J$-module along an automorphism $\alpha$ of a free
$V/J$-module. In general $\alpha$ need not lift to $V$. But $\operatorname{diag}(\alpha,\alpha^{-1})$ is a product
of elementary matrices, and these always lift. Hence $\Int(R)$ is a direct summand of a free module. Since $R$ is
local, Kaplansky's theorem makes it free.

Where the hypotheses are used:
\begin{itemize}
\item That $V$ is a discrete valuation ring with finite residue field gives the Newton basis (Lemma~\ref{lem:newton})
and the digit polynomial (Lemma~\ref{lem:indicators}).
\item The conductor $\pi^mV\subseteq R$ gives the square and the finiteness of $V/J$.
\item Locality of $R$ gives freeness from projectivity.
\end{itemize}
Corollary~\ref{cor:ai} reduces to the theorem: for an analytically irreducible domain the normalization is a finite,
local discrete valuation ring.

\section{The ring and the module over the valuation ring}

From now on $V,\pi,K,m$ are as in Theorem~\ref{thm:main}. Let $Q$ be the order of the residue field of $V$, and $v$
the normalized valuation. Put $J=\pi^mV$, $\bar D=V/J$ (a finite ring), $R_0=R/J\subseteq\bar D$,
\[
 A=\Int(R),\qquad B=\Int(R,V)=\{f\in K[X]:f(R)\subseteq V\}.
\]
Then $A\subseteq B$, because $R\subseteq V$.

\begin{lemma}\label{lem:ring}
$R$ is a one-dimensional Noetherian local domain with fraction field $K$ and maximal ideal $R\cap\pi V$. Moreover
$V$ is a finite $R$-module and is the integral closure of $R$.
\end{lemma}

\begin{proof}
Every $v\in V$ equals $(\pi^mv)/\pi^m$ with both terms in $R$, so $\operatorname{Frac}(R)=K$.

\emph{Local.} Let $x\in R$ with $v(x)=0$. Multiplication by $x$ on the finite ring $R_0$ is injective, since $x$ is a
unit of $V$ and $J$ is an ideal of $V$. Hence it is surjective, and there is $y\in R$ with $xy\in1+J$. The inverse of
an element of $1+J$ in $V$ lies in $1+J\subseteq R$, so $x$ is a unit of $R$.

\emph{Finite.} $V$ is generated over $R$ by representatives of the finite set $V/J$.

\emph{Noetherian.} Let $H\neq0$ be an ideal and $h\in H$ of least valuation. Then $JhV=hJ\subseteq H\subseteq hV$,
and $hV/hJ$ is finite. The ideal $hJ=h\pi^mV$ is isomorphic to $V$ as an $R$-module, so it is finitely generated. Finitely many further elements of $H$ represent its image in
$hV/hJ$. So $H$ is finitely generated.

\emph{Dimension.} Every nonzero prime contains some $hJ$, so its residue ring is finite, hence a field.

\emph{Integral closure.} $V$ is finite, hence integral, over $R$, and it is integrally closed in $K$.
\end{proof}

\begin{lemma}\label{lem:newton}
Choose $a_0\in R$. Inductively choose $a_n\in R$ minimizing $v\bigl(\prod_{i<n}(a_n-a_i)\bigr)$, and put
$b_n=\prod_{i<n}(X-a_i)/(a_n-a_i)$. Then $B=\bigoplus_{n\geq0}Vb_n$.
\end{lemma}

\begin{proof}
$R$ is infinite, since it contains $J$, so the minimum is finite and the $a_n$ are distinct. By minimality,
$v(\prod_{i<n}(x-a_i))\geq v(\prod_{i<n}(a_n-a_i))$ for every $x\in R$, so $b_n\in B$; also $b_n(a_n)=1$ and
$b_n(a_i)=0$ for $i<n$. Let $f\in B$ have degree $d$, and write $f=\sum_{n\le d}c_nb_n$ with $c_n\in K$. Evaluating
at $a_0,a_1,\dots$ in turn shows that $c_0=f(a_0)\in V$, then $c_1=f(a_1)-c_0b_0(a_1)\in V$, and so on: each
$c_n=f(a_n)-\sum_{i<n}c_ib_i(a_n)\in V$.
\end{proof}

This is Bhargava's construction of a regular basis from a $\pi$-ordering \cite{Bha}.

\section{Reduction modulo the conductor}

For $N\geq1$ let $E_N$ be the set of classes of $V/\pi^NV$ that meet $R$. For a ring $S$, let $\Fn(R,S)$ be the ring
of functions $R\to S$ that factor through $R\to E_N$ for some $N$. It is the union of the rings $S^{E_N}$, embedded
by pullback along $E_{N+1}\to E_N$.

\begin{lemma}\label{lem:indicators}
Evaluation induces ring isomorphisms
\[
 B/JB\xrightarrow{\ \sim\ }\Fn(R,\bar D),\qquad A/JB\xrightarrow{\ \sim\ }\Fn(R,R_0).
\]
Note that $JB\subseteq A$.
\end{lemma}

\begin{proof}
If $f\in JB$ then $f(R)\subseteq J\subseteq R$, so $JB\subseteq A$.

\emph{Well defined.} If the coefficients of $f\in B$ have valuation at least $-h$, then $x-y$ divides $x^k-y^k$ in
$V$, so $v(f(x)-f(y))\geq v(x-y)-h$. Hence $f\bmod J$ factors through $E_{m+h}$.

\emph{Injective.} If $f(R)\subseteq J$, then $\pi^{-m}f\in B$, so $f\in JB$.

\emph{Surjective.} Let $\delta=(X^Q-X)/\pi$. It maps $V$ to $V$, because $x^Q\equiv x\pmod\pi$. If $v(x-y)=r\geq1$,
then
\[
 \delta(x)-\delta(y)=\frac{x-y}{\pi}\Bigl(\sum_{i=0}^{Q-1}x^{Q-1-i}y^i-1\Bigr),
\]
and the bracket is $\equiv Qx^{Q-1}-1\equiv-1\pmod\pi$, since $Q$ is a power of the residue characteristic. So
$v(\delta(x)-\delta(y))=r-1$. By induction, for $x,y\in V$,
\[
 x\equiv y\pmod{\pi^N}\iff \delta^{\circ j}(x)\equiv\delta^{\circ j}(y)\pmod\pi\ \text{ for } 0\leq j<N .
\]
Indeed, if $v(x-y)=r<N$, the valuations of $\delta^{\circ j}(x)-\delta^{\circ j}(y)$ are $r,r-1,\dots,0$ for
$j=0,\dots,r$. If $r\geq N$, they are at least $1$ for $j<N$.

Fix a class $c\in E_N$, a point $x_c\in R$ of it, and $c_j=\delta^{\circ j}(x_c)$. Put
$H_c=\prod_{j<N}\bigl(1-(\delta^{\circ j}(X)-c_j)^{Q-1}\bigr)$. Then $H_c$ maps $V$ to $V$, and for $x\in V$:
\begin{itemize}
\item if $x$ lies in $c$, every factor is $\equiv1\pmod\pi$, so $H_c(x)\equiv1$;
\item otherwise some factor is $1-u^{Q-1}$ with $u$ a unit, which is $\equiv0\pmod\pi$, so $H_c(x)\equiv0$.
\end{itemize}
Raise $H_c$ to the power $Q^m$:
\begin{itemize}
\item if $z\in\pi V$, then $v(z^{Q^m})\geq Q^m\geq m$;
\item if $z=1+w$ with $v(w)\geq1$, then $v((1+w)^p-1)\geq v(w)+1$, because the middle binomial coefficients are
divisible by $p$ and $v(w^p)\geq v(w)+1$. After the $\log_pQ^m\geq m$ successive $p$-th powers,
$z^{Q^m}\equiv1\pmod{\pi^m}$.
\end{itemize}
So $H_c^{Q^m}\in B$ is the indicator of $c$ modulo $J$. A function in $\Fn(R,\bar D)$ that factors through
$E_N$ is $\sum_{c\in E_N}d_cH_c^{Q^m}\bmod J$, with $d_c\in V$ lifting its values.

\emph{The image of $A$.} Values of $f\in A$ lie in $R$, so they reduce into $R_0$. Conversely, if $f\in B$ has all
its values in $R_0$ modulo $J$, then $f(R)\subseteq R+J=R$.
\end{proof}

\begin{lemma}\label{lem:free-functions}
$C=\Fn(R,R_0)$ is a free $R_0$-module of countably infinite rank, and $C\otimes_{R_0}\bar D\to\Fn(R,\bar D)$ is an
isomorphism.
\end{lemma}

\begin{proof}
Choose for each $N$ a map $s_N:E_N\to E_{N+1}$ that picks one class above each class. Restriction along $s_N$ is a
retraction of $R_0^{E_N}\to R_0^{E_{N+1}}$, and its kernel is free on $E_{N+1}\setminus s_N(E_N)$. Hence
$C=R_0^{E_1}\oplus\bigoplus_N\ker$ is free. Its rank is infinite, because $R$ is infinite and distinct points lie in
distinct classes for large $N$. The tensor product commutes with the union, and $R_0^{E_N}\otimes_{R_0}\bar D=\bar
D^{E_N}$.
\end{proof}

\section{Patching, and the proof of Theorem~\ref{thm:main}}

\begin{lemma}\label{lem:patch}
Let $V\to\bar D$ be a surjective ring map, $R_0\to\bar D$ a ring map, and $R=V\times_{\bar D}R_0$. Let $I$ be a set
and $\alpha$ an automorphism of $\bar D^{(I)}$. Put
\[
 M_\alpha=\{(b,c)\in V^{(I)}\oplus R_0^{(I)}:\ \bar b=\alpha(\bar c)\}.
\]
Then $M_\alpha\oplus M_{\alpha^{-1}}\cong R^{(I\sqcup I)}$. In particular $M_\alpha$ is projective over $R$.
\end{lemma}

\begin{proof}
Matrices act on direct sums, so their columns have finite support. In $\operatorname{End}(\bar D^{(I)}\oplus\bar
D^{(I)})$, block multiplication gives
\[
 \begin{pmatrix}\alpha&0\\0&\alpha^{-1}\end{pmatrix}
 =\begin{pmatrix}1&\alpha\\0&1\end{pmatrix}
 \begin{pmatrix}1&0\\-\alpha^{-1}&1\end{pmatrix}
 \begin{pmatrix}1&\alpha\\0&1\end{pmatrix}
 \begin{pmatrix}0&-1\\1&0\end{pmatrix}.
\]
Lift each entry of $\alpha$ and $\alpha^{-1}$ to $V$, lifting zero to zero, so that columns keep finite support.
The four factors with these lifts are automorphisms of $V^{(I\sqcup I)}$, since $\begin{pmatrix}1&a\\0&1\end{pmatrix}$
has inverse $\begin{pmatrix}1&-a\\0&1\end{pmatrix}$ whatever $a$ is. Their product $\Gamma$ reduces to
$\operatorname{diag}(\alpha,\alpha^{-1})$. Now $(b,c)\mapsto(\Gamma^{-1}b,c)$ maps $M_\alpha\oplus M_{\alpha^{-1}}$
isomorphically onto $\{(b,c):\bar b=\bar c\}=R^{(I\sqcup I)}$.
\end{proof}

\begin{proof}[Proof of Theorem~\ref{thm:main}]
Lemma~\ref{lem:ring} gives the properties of $R$. Since $J\subseteq R$, the map $R\to V\times R_0$ identifies $R$
with $V\times_{\bar D}R_0$.

By Lemma~\ref{lem:newton}, $B\cong V^{(I)}$ with $I$ countably infinite. By Lemmas~\ref{lem:indicators}
and~\ref{lem:free-functions}, $C=A/JB\cong R_0^{(I)}$, and the composite
\[
 C\otimes_{R_0}\bar D\cong\Fn(R,\bar D)\cong B/JB
\]
is an isomorphism. In these bases it is an automorphism $\alpha$ of $\bar D^{(I)}$. An element of $A$ is exactly an
$f\in B$ whose class modulo $JB$ lies in $C$ (Lemma~\ref{lem:indicators}). So $f\mapsto(f,f\bmod JB)$ identifies $A$
with $M_\alpha$, which is projective by Lemma~\ref{lem:patch}.

A projective module over a local ring is free \cite[Theorem~2]{Kap}. $A$ is countably generated, and it is not
finitely generated, since it contains $R[X]$ and is torsion-free of infinite rank. So it is free of countably
infinite rank.
\end{proof}

\section{Analytically irreducible domains}

\begin{proof}[Proof of Corollary~\ref{cor:ai}]
Let $(D,M)$ be as in the corollary, with fraction field $K$.

\emph{Infinite residue field.} Let $f\in\Int(D)$ have degree $n$. Choose $a_0,\dots,a_n\in D$ with distinct residues.
Lagrange interpolation gives
\[
 f=\sum_i f(a_i)\prod_{j\neq i}\frac{X-a_j}{a_i-a_j},
\]
whose denominators are units. Hence $\Int(D)=D[X]$, which is free.

\emph{Finite residue field.} Since the completion $\hat D$ is a domain, it is reduced, so $D$ is analytically
unramified. Its integral closure $D'$ in $K$ is therefore a finite $D$-module \cite[Tag~032Y]{Stacks}. So $D'$ is a
one-dimensional integrally closed Noetherian domain, that is, a semilocal Dedekind domain. Each of its maximal
ideals $\mathfrak n_1,\dots,\mathfrak n_r$ lies over $M$.

We show $r=1$. By \cite[Proposition~10.13]{AM}, $D'\otimes_D\hat D$ is the inverse limit of the rings $D'/M^kD'$.
Each of these has finite length, so it is the product of its localizations at the $\mathfrak n_i$, all nonzero
\cite[Theorem~8.7]{AM}. In the limit, $D'\otimes_D\hat D$ is a product of $r$ nonzero rings.

On the other hand, $\hat D$ is flat over $D$ \cite[Proposition~10.14]{AM}. So $D'\otimes_D\hat D\subseteq
K\otimes_D\hat D=S^{-1}\hat D$, where $S=D\setminus\{0\}$, and the elements of $S$ are nonzerodivisors on $\hat D$.
Thus $D'\otimes_D\hat D$ lies in the fraction field of the domain $\hat D$ and has no nontrivial idempotents.
So $r=1$.

Hence $V=D'$ is a discrete valuation ring with finite residue field, finite over $D$. Take a nonzero $s\in D$ with
$sD'\subseteq D$, which exists because $D'$ is finite. Then the conductor $\{x\in D':xD'\subseteq D\}$ is a nonzero
ideal of $V$, so it contains $\pi^mV$ for some $m\geq1$. Theorem~\ref{thm:main} applies.

Elliott's two rings are complete, hence analytically irreducible. This refutes~\eqref{eq:elliott}, and with it
Conjectures~2.15 and~2.17. Conjecture~2.10 asserts the existence of a domain that the corollary excludes.
\end{proof}

\section{Elliott's two rings by hand}

The following identities give the failure of~\eqref{eq:elliott} without the general theorem. They were also checked
by exact series arithmetic, with the program \texttt{check\_examples.py}, which is kept with the records of this note.

\begin{example}\label{ex:f4}
Let $V=\F_4[[t]]$ with $\F_4=\F_2(\omega)$ and $\omega^2=\omega+1$, let $D=\F_2+tV$, and let $M=tV$. Every $x\in D$
is $a+tv$ with $a\in\F_2$ and $v\in V$. Since $a^2=a$,
\[
 q(x)=\frac{x^2+x}{t}=v+tv^2\in V .
\]
Write the constant term of $q(x)$ as $\alpha+\beta\omega$ with $\alpha,\beta\in\F_2$. Then $q(x)^2$ has constant term
$\alpha+\beta\omega^2=\alpha+\beta+\beta\omega$, so:
\begin{itemize}
\item $h_2=q+q^2$ takes the value $\beta+tV\subseteq D$;
\item $h_1=q+\omega h_2$ takes the value $\alpha+tV\subseteq D$.
\end{itemize}
Thus $h_1,h_2\in\Int(D)$. In characteristic $2$, $q=h_1+\omega h_2$, so
\[
 X^2+X=t\,h_1+(\omega t)\,h_2\in M\,\Int(D).
\]
Here $q\in\Int(D,V)$ but $q\notin\Int(D)$: for instance $q(\omega t)=\omega+\omega^2t$.
\end{example}

\begin{example}\label{ex:cusp}
Let $V=\F_2[[t]]$, $D=\F_2[[t^2,t^3]]=\F_2+t^2V$, and $M=(t^2,t^3)=t^2V$. Every $x\in D$ is $a+t^2v$, and
\[
 q(x)=\frac{x^2+x}{t^2}=v+t^2v^2=v_0+v_1t+O(t^2).
\]
Since $q(x)^2=v_0+O(t^2)$, the element $y=(q(x)+q(x)^2)/t$ lies in $V$ and satisfies $y\equiv v_1\pmod t$. Frobenius
gives $y^4\in\F_2+t^4V\subseteq D$ and $y^4\equiv v_1\pmod{t^4}$. So:
\begin{itemize}
\item $h_2=\bigl((q+q^2)/t\bigr)^4$ lies in $\Int(D)$;
\item $h_1=q+t\,h_2$ takes the value $v_0+v_1t+t\,v_1+O(t^2)=v_0+O(t^2)$, which lies in $D$.
\end{itemize}
Thus $h_1\in\Int(D)$, and
\[
 X^2+X=t^2h_1+t^3h_2\in M\,\Int(D).
\]
Again $q\notin\Int(D)$: the value $q(t^3)=t+t^4$ has a nonzero coefficient of $t$.
\end{example}

In both examples, the polynomial that Elliott proposed as a witness of non-flatness, $(X^2+X)/t^m$, lies in
$V\Int(D)$, as Theorem~\ref{thm:main} predicts. In Ma's two-branch domain the analogous polynomial does not
\cite[Proposition~4.1]{Ma}.

\section{What remains open}

\begin{itemize}
\item \emph{Infinite normalization.} Theorem~\ref{thm:main} needs a finite normalization. There are one-dimensional
Noetherian local domains with finite residue field whose integral closure is local but not finite (analytically
ramified unibranch domains). Whether $\Int(D)$ is flat or free for them is open. The proof uses the conductor, which
is zero there.
\item \emph{Several branches.} When the finite normalization has several maximal ideals, Ma's example shows that
flatness can fail. We do not characterize which multi-branch domains have flat $\Int(D)$.
\item \emph{Higher dimension.} We do not treat local domains of dimension at least two outside the Krull and TV PvMD
cases.
\item \emph{Regular bases.} The basis of $\Int(D)$ produced here is not shown to be regular, meaning one polynomial
of each degree, and we do not decide whether a regular basis exists for Elliott's rings.
\end{itemize}

\section*{Note on method}
The proof of Theorem~\ref{thm:main} was produced by GPT-6 Astra (ChatGPT Pro mode), in a study of the operations on
integer-valued polynomials that arose from Grothendieck's Montpellier manuscripts. The study also checked the proof
in a targeted reading.

Claude Opus 5.5 at effort high checked every step of that proof, and wrote this text. It also wrote the reduction
of Corollary~\ref{cor:ai}, the reading of Corollary~\ref{cor:biring} through Elliott's theorems, and the two
examples. The examples were verified by an exact program. Claude Opus 5.5 at effort max then checked every step of the
note again, the corollary and the examples included.

The sources \cite{CFFG,E11a,E11b,E16b,Ma,JLSWXY} were read at the locators given, and \cite[Tag~032Y]{Stacks} at
the source. The results of \cite{E07,E10} are cited as \cite[pp.~9--10]{CFFG} and \cite[p.~16]{E11b} recall them;
we have not opened \cite{E07,E10}. The monograph \cite{CC} was not consulted.

No human has checked this note. Our text is dedicated to the public domain under CC0~1.0. The program
\texttt{check\_examples.py} and the records of the checks are at the permanent page of this work, WB-0004:
\url{\workbenchurl}.

\begin{thebibliography}{99}\small
\bibitem{AM} M.~F. Atiyah and I.~G. Macdonald, \emph{Introduction to Commutative Algebra}, Addison-Wesley, 1969.
\bibitem{Bha} M.~Bhargava, P-orderings and polynomial functions on arbitrary subsets of Dedekind rings,
\emph{J. reine angew. Math.} 490 (1997), 101--127.
\bibitem{CC} P.-J. Cahen and J.-L. Chabert, \emph{Integer-Valued Polynomials}, Math. Surveys Monogr.~48, Amer. Math.
Soc., 1997.
\bibitem{CFFG} P.-J. Cahen, M. Fontana, S. Frisch and S. Glaz, Open problems in commutative ring theory, in
\emph{Commutative Algebra} (M.~Fontana, S.~Frisch, S.~Glaz, eds.), Springer, 2014, 353--375; page numbers refer to
the authors' version of 23 December 2013.
\bibitem{E07} J.~Elliott, Universal properties of integer-valued polynomial rings, \emph{J. Algebra} 318 (2007),
68--92.
\bibitem{E10} J.~Elliott, Birings and plethories of integer-valued polynomials, \emph{Actes des Rencontres du CIRM}
2 (2010), no.~2.
\bibitem{E11a} J.~Elliott, Integer-valued polynomial rings, $t$-closure, and associated primes, \emph{Comm.
Algebra} 39 (2011), 4128--4147; page numbers refer to arXiv:1105.0142v1.
\bibitem{E11b} J.~Elliott, Presentations and module bases of integer-valued polynomial rings, arXiv:1109.3921v1
(2011).
\bibitem{E16b} J.~Elliott, Integer-valued polynomials on commutative rings and modules, arXiv:1608.00171v1 (2016).
\bibitem{JLSWXY} J.~Jiang, Y.~Li, Z.~Sun, Y.~Wang, L.~Xiao and J.~Yu, On some open problems in commutative algebra
resolved by Rethlas, arXiv:2605.25259v2 (2026).
\bibitem{Rethlas} H.~Ju et al., Automated conjecture resolution with formal verification, arXiv:2604.03789 (2026).
\bibitem{Kap} I.~Kaplansky, Projective modules, \emph{Ann. of Math.} 68 (1958), 372--377.
\bibitem{Ma} H.~Ma, A counterexample to Problem 19 on integer-valued polynomial rings, arXiv:2604.05922v2 (2026).
\bibitem{Mil} J.~Milnor, \emph{Introduction to Algebraic K-Theory}, Ann. of Math. Stud.~72, Princeton, 1971.
\bibitem{Stacks} The Stacks Project Authors, \emph{The Stacks Project}, \url{https://stacks.math.columbia.edu}.
\end{thebibliography}
\end{document}
